
% Module_8A_cmb_tensor_test.tex
\documentclass[12pt]{article}
\usepackage{amsmath,amssymb,hyperref,booktabs}
\usepackage{geometry}
\geometry{margin=1in}
\title{Module 8-A: Imminent CMB Probe of the Tensor Prediction}
\author{}
\date{}
\begin{document}
\maketitle

\section*{Motivation}
The SAT filament‑inflation module predicts a \emph{tensor‑to‑scalar ratio}
\[
  r_{0.05} \;=\; 0.14 \;\pm\; 0.02 \quad (95\%\,\text{CL})
\tag{8A.1}
\]
at pivot scale $k=0.05\;\text{Mpc}^{-1}$, set entirely by the
\emph{torsion‑inflaton coupling} $g_\tau = 3.7$ fixed in Module~8.  
Current CMB limits already push below this value, so the theory faces a decisive
test within the next decade.

\section*{O8‑A.1 \quad Current Status (2025)}
\begin{itemize}
  \item BK18 + Planck 2018 + WMAP: $r_{0.05} < 0.036$ (95\% CL).
  \item Simons Observatory early science (forecast): $\sigma(r)\simeq 0.009$ by 2029.
\end{itemize}
SAT thus sits $\approx12\sigma$ \emph{above} the present upper bound—an uncomfortable but falsifiable tension.

\section*{O8‑A.2 \quad LiteBIRD Reach}
LiteBIRD’s baseline design (1.4 m telescope, 34 GHz–402 GHz bands, full‑sky, $\ell=2$–400) yields
\[
  \sigma(r) \;\simeq\; 0.0025 \quad\Longrightarrow\quad
  2\sigma\; \text{reach}: r_{0.05}\approx 0.005.
\tag{8A.2}
\]
Even under pessimistic foreground residuals (30\% synchrotron, 10\% dust), 
$\sigma(r)<0.004$ remains achievable, far below Eq.~\eqref{8A.1}.  

\begin{center}
\renewcommand{\arraystretch}{1.15}
\begin{tabular}{@{}lccc@{}}
\toprule
Experiment & Year & Band coverage & $\sigma(r_{0.05})$ \\
\midrule
BK18 (final) & 2025 & 95–270 GHz & -- $(<$0.012)$ \\
Simons Obs.  & 2029 & 27–280 GHz & 0.009 \\
\textbf{LiteBIRD} & 2033 & 34–402 GHz & \textbf{0.0025} \\
CMB‑S4 & 2035 & 20–270 GHz & 0.001 \\
\bottomrule
\end{tabular}
\end{center}

\paragraph{Implication.}
If the SAT prediction stands, LiteBIRD should detect $r$ at
\[
  \frac{0.14}{0.0025} \;\approx\; 56\sigma,
\]
well beyond discovery threshold.

\section*{O8‑A.3 \quad Contingency: What If $r<0.03$?}
If LiteBIRD (or Simons Obs.) reports $r_{0.05}<0.03$ (95\% CL),
the torsion‑inflaton sector in its current form is ruled out.
Possible salvage options:
\begin{enumerate}
  \item \emph{Mode 2 correction}: allow $\Sigma_t$ triggered dynamics to
        \textbf{dilute} primordial torsion by a factor
        $e^{-\beta N_e}$ during the last $N_e\simeq50$
        $e$‑folds.  Requires $\beta\simeq0.06$—in tension with
        renormalisability unless a new scale appears.
  \item \emph{Hybrid inflation extension}: embed a spectator scalar
        $\chi$ coupled $g_{\chi\tau}\chi^2\tau^2$; tuning $g_{\chi\tau}\lesssim10^{-3}$ can
        lower $r$ while respecting Mode 1 priors, but adds a free parameter.
\end{enumerate}
Both routes violate the ``no extra knobs'' philosophy; we therefore
advertise the $r$ test as a \emph{clean kill} for the minimalist SAT.

\section*{O8‑A.4 \quad Associated Observables}
The same fit yields 
\[
  n_s = 0.957\;(\pm0.004), \qquad
  \alpha_s = -0.0035\;(\pm0.0011),
\]
consistent with Planck but correlated with $r$ via
$n_s \approx 1 - \frac{r}{4}$.  
Detecting $r\simeq0.14$ implies $n_s\simeq0.965$, a value LiteBIRD
will measure to $\sigma(n_s)\approx0.002$, offering an internal
consistency cross‑check.

\section*{Key Takeaways}
\begin{itemize}
  \item The minimalist SAT predicts $r=0.14$—already high relative to 2025 bounds.
  \item LiteBIRD (launch 2030, first data 2033) will reach $\sigma(r)\approx0.0025$,
        guaranteeing either a $>50\sigma$ detection or a decisive falsification.
  \item A null result ($r<0.03$) would force an \emph{extra‑parameter} patch, undercutting the theory’s minimalism.
\end{itemize}

\end{document}
